\section{Related Work}

\paragraph{Benchmarks for Conversational AI Agents.}
Following a long line of research into language agents~\cite{yao2022webshop, zhou2023webarena, jimenez2023swe, liu2023agentbench, ruan2023identifying}, LLM tool use~\cite{berkeley-function-calling-leaderboard,qin2024toolllm,huang2023metatool}, and task-oriented dialog~\cite{chen2021action, budzianowski2018multiwoz, andreas2020task,schatzmann2007evaluating,gur2018hidden,he2018decoupling,Hu_2023}, \taubench~\cite{yao2024tau} is a recently introduced benchmark to measure the reliability of language agents in multi-turn task-oriented conversations such as customer service workflows, while respecting domain rules. Each task in \taubench instantiates a live conversation between a user simulator and the language agent, with tasks spread across two domains -- \retail and \airline. To quantify reliability the paper introduces \passk metrics: the fraction of \textit{k} independent runs that succeed.

Several follow-ups to \taubench have explored variations of the basic setup. FlowBench~\cite{xiao2024flowbench} isolates the planning step of tool-using agents by injecting explicit workflow knowledge into the prompt, using natural language, python-like pseudocode or mermaid flowcharts. IntellAgent~\cite{levi2025intellagent} provides an evaluation pipeline to programmatically build synthetic test suites from structured policy graphs that encode domain rules and their co-occurrence statistics. IntellAgent explicitly uses \taubench as an external gold standard, reporting a high Spearman correlation between the two score distributions, and acts as a fast, synthetic proxy task. APIGen-MT~\cite{prabhakar2025apigen} explores the idea of fine-tuning tool-calling agents for \taubench. They generate data by creating conversation blueprints which are sequences of tool calls that depend on each other, followed by simulating conversational traces based on each blueprint. ToolSandbox~\cite{lu2024toolsandbox} focuses on creating stateful tools in order to evaluate agent progress in a more fine-grained manner.

Our work extends the \taubench paradigm and generalizes it to allow for both the user and agent to have state-changing abilities (via tool calls) over a shared world. As demonstrated in the results, this allows us to build more complex domains to test conversational agents, while also providing the opportunity for fine-grained analysis of agent failure points that can be improved upon.

\paragraph{User Simulation for Conversational Agents.}
The reliability of user simulation has been a key concern for benchmarks like \taubench~\cite{yao2024tau}. While most efforts have focused on introducing supervision for the user simulator, for instance, by using a generic LLM to generate or validate user responses~\cite{prabhakar2025apigen}, less attention has been paid to the possibility of using the environment itself to constrain and shape user simulator behavior for increased reliability, a core tenet of our approach. This concern has been extensively studied in the context of task-oriented dialogue systems, with early work by \cite{pietquin_survey_2013} providing a comprehensive survey of metrics for evaluating user simulations. More recently, \cite{kazi_large_2024} has demonstrated how LLMs can be effectively used as user-agents for evaluating task-oriented dialogue systems, showing that careful prompting and state tracking can lead to more reliable and context-aware user simulations.

\paragraph{Multi-Agent Benchmarks.}
Our work is also related to efforts to build multi-agent frameworks and evaluate them~\cite{zhu2025multiagentbench}.
While we can consider the user and the agent in our paper as forming a multi-agent system, the key difference in our case is that the final evaluation still focuses on the agent's ability to elicit the right information from the user and perform the correct actions to solve the task. This introduces an inherent asymmetry between the agent and the user; our focus is not on solving a pure multi-agent problem but rather on the agent's capability to effectively guide and collaborate with a user who also possesses agency. In this sense, the framework can be collaborative (e.g., troubleshooting), competitive (e.g., negotiating a subscription), or a hybrid, requiring the agent to identify and navigate the scenario appropriately, even accounting for user mistakes or errors.
